\documentclass[12pt, letterpaper]{article}

% ----- CONFIGURACIÓN APA 7MA EDICIÓN -----

% Márgenes de 2.54 cm (1 pulgada) en los cuatro bordes
\usepackage[margin=1in, headheight=1.5cm, headsep=0.5cm]{geometry}

% Fuente Times New Roman
\usepackage{mathptmx}

% Alineación a la izquierda (texto sin justificar) y sin guiones al final
\usepackage[document]{ragged2e}
\usepackage[none]{hyphenat}

% Interlineado doble
\usepackage{setspace}
\doublespacing

% Sangría de 1.27 cm (0.5 pulgadas) en la primera línea de cada párrafo
\setlength{\parindent}{0.5in}

% Paginación y encabezados personalizados
\usepackage{fancyhdr}
\pagestyle{fancy}
\fancyhf{} % Limpiar encabezados y pies
\renewcommand{\headrulewidth}{1pt}
\renewcommand{\footrulewidth}{1pt}
\fancyhead[L]{
\includegraphics[height=1.2cm]{IMAGENES/CARATULA.png}}
\fancyhead[R]{\small\textbf{ESCUELA DE EDUCACIÓN SUPERIOR TECNOLÓGICA PRIVADA LA}\\\textbf{DIRECCIÓN ACADÉMICA}}
\fancyfoot[L]{\small Informe Ejecutivo de Experiencias Formativas en Situaciones Reales de Trabajo}
\fancyfoot[R]{\small pág. \thepage}


% ----- CONFIGURACIÓN DE TÍTULOS (JERARQUÍA APA 7) -----
\usepackage{titlesec}
\setcounter{secnumdepth}{0} % Sin numeración de títulos

% Nivel 1: Centrado, Negrita
\titleformat{\section}[block]{\normalfont\bfseries\centering}{\thesection}{0pt}{}
\titlespacing*{\section}{0pt}{\baselineskip}{0pt}

% Nivel 2: Alineado a la Izquierda, Negrita
\titleformat{\subsection}[block]{\normalfont\bfseries\raggedright}{\thesubsection}{0pt}{}
\titlespacing*{\subsection}{0pt}{\baselineskip}{0pt}

% Nivel 3: Alineado a la Izquierda, Negrita y Cursiva
\titleformat{\subsubsection}[block]{\normalfont\bfseries\itshape\raggedright}{\thesubsubsection}{0pt}{}
\titlespacing*{\subsubsection}{0pt}{\baselineskip}{0pt}

% Nivel 4: Sangría de 1.27 cm, Negrita, finalizado con punto (runin)
\titleformat{\paragraph}[runin]{\normalfont\bfseries}{\theparagraph}{0pt}{}[.]
\titlespacing*{\paragraph}{\parindent}{\baselineskip}{0.5em}

% Nivel 5: Sangría de 1.27 cm, Negrita y Cursiva, finalizado con punto (runin)
\titleformat{\subparagraph}[runin]{\normalfont\bfseries\itshape}{\thesubparagraph}{0pt}{}[.]
\titlespacing*{\subparagraph}{\parindent}{\baselineskip}{0.5em}


% ----- CONFIGURACIÓN DE TABLAS Y FIGURAS -----
\usepackage{booktabs}
\usepackage{caption}
\usepackage{graphicx}
\usepackage{float}
\usepackage{array}

% Formato para títulos de tablas y figuras según APA 7
\captionsetup[table]{labelsep=newline, justification=raggedright, singlelinecheck=false, textfont=it, labelfont=bf, skip=0pt}
\captionsetup[figure]{labelsep=newline, justification=raggedright, singlelinecheck=false, textfont=it, labelfont=bf, skip=0pt}

% ----- IDIOMA Y ENCODING -----
\usepackage[utf8]{inputenc}
\usepackage[spanish]{babel}

% Paquete para el índice con el formato adecuado
\usepackage{tocloft}
\renewcommand{\cftsecleader}{\cftdotfill{\cftdotsep}}


% ----- ENLACES CLICKABLES (NDICES REDIRECCIONABLES) -----
\usepackage[colorlinks=true, linkcolor=black, urlcolor=blue, citecolor=black, bookmarks=true]{hyperref}

\begin{document}
\RaggedRight % Aplicar texto sin justificar

% ----- CARÁTULA -----
\begin{titlepage}
    \thispagestyle{fancy} % Forzar encabezado y pie en la carátula
    \begin{center}
        \singlespacing 
        
        \vspace*{0.5cm}
        \textbf{ESCUELA DE EDUCACIÓN SUPERIOR TECNOLÓGICA PRIVADA LA PONTIFICIA} \\
        \vspace{0.5cm}
        
        Informe Ejecutivo de Experiencias Formativas en Situaciones Reales de Trabajo IV \\
        en la modalidad de Proyecto Productivo y/o Proyecto de Innovación \\
        \vspace{1.5cm}
        
        % Logo principal de la carátula
        
\includegraphics[width=8cm]{IMAGENES/CARATULA.png} \\
        \vspace{1.5cm}
        
        \textbf{Informe Ejecutivo de Experiencias Formativas en Situaciones Reales de Trabajo IV} \\
        \textbf{de la carrera profesional de Ingeniería de Sistemas de Información} \\
        \vspace{1.5cm}
        
        Informe correspondiente al Módulo Formativo IV denominado \\
        \vspace{0.5cm}
        \textbf{IMPLEMENTACIÓN DE PRUEBAS DE SOFTWARE} \\
        \vspace{0.5cm}
        ParkSys — Sistema de Gestión de Estacionamiento Vehicular para el establecimiento \\
        ParkSys Central Ayacucho. \\
        \vspace{1cm}
        
        \textbf{Docente:} \\
        ING. SAUÑE PICHARDO, JHONATAN \\
        \vspace{1.5cm}
        
        \textbf{desarrollado por el estudiante} \\
        ASPUR YARANGA, JHOEL EMILIO \\
        CONDORI COTAQUISPE, SULY JHOMY \\
        \vspace{1.5cm}
        
        Huamanga, 04 de septiembre del 2026
        
    \end{center}
\end{titlepage}

\doublespacing

% ----- ÍNDICES -----
\newpage
\tableofcontents
\newpage
\listoffigures
\newpage
\listoftables

% ----- INCLUSIÓN DE CAPÍTULOS -----
\clearpage
\section{INTRODUCCION}

El presente documento constituye el Informe Ejecutivo de Experiencias Formativas en Situaciones Reales de Trabajo IV, en la modalidad de Proyecto Productivo y/o Proyecto de Innovación, desarrollado para la carrera profesional de Ingeniería de Sistemas de Información de la Escuela de Educación Superior Tecnológica Privada La Pontificia. 

En este informe se detalla el desarrollo, implementación, pruebas y despliegue del sistema ParkSys, una plataforma integral de gestión inteligente de estacionamientos (Smart Parking) concebida para automatizar y optimizar el ciclo de vida del proceso de aparcamiento en el establecimiento ParkSys Central Ayacucho. El informe está estructurado en tres capítulos principales: el Capítulo 1 aborda el desarrollo e implementación del sistema, incluyendo la arquitectura, base de datos y tecnologías utilizadas; el Capítulo 2 detalla las pruebas de software, optimización de rendimiento y el proceso de despliegue continuo; finalmente, el Capítulo 3 presenta las conclusiones derivadas de la experiencia y las recomendaciones pertinentes.


\clearpage
\section{CAPÍTULO 1: DESARROLLO E IMPLEMENTACIÓN DEL SISTEMA}

\subsection{1.1. DESCRIPCIÓN DEL SISTEMA ELEGIDO}

El sistema desarrollado como parte de las presentes Experiencias Formativas en Situaciones Reales de Trabajo IV es ParkSys, una plataforma integral de gestión inteligente de estacionamientos (\textit{Smart Parking}), concebida para automatizar y optimizar de forma completa el ciclo de vida del proceso de aparcamiento, desde la reserva anticipada hasta el cobro final por el uso del espacio físico.

\subsubsection{Propósito del sistema}

ParkSys busca resolver los principales cuellos de botella operativos que enfrentan los establecimientos de estacionamiento tradicionales, tales como el control manual del acceso vehicular, el fraude o error humano en los cobros, y la falta de visibilidad en tiempo real sobre la disponibilidad de espacios. El sistema ofrece dos modalidades de uso simultáneas:

\begin{enumerate}
    \item \textbf{Registro in-situ (Walk-in):} controlado directamente por los operadores de garita física cuando un vehículo arriba sin reserva previa.
    \item \textbf{Reserva anticipada (Self-service):} a través de un portal digital responsivo dirigido al conductor, con selección geométrica de cajón en mapa 2D y pasarela de pago en línea.
\end{enumerate}

\subsubsection{Módulos principales del sistema}

\begin{table}[hbt!]
\caption{Módulos principales del sistema}
\label{tab:modulos}
\begin{tabular}{@{}p{0.5cm} p{3.5cm} p{11cm}@{}}
\toprule
\textbf{N.º} & \textbf{Módulo} & \textbf{Descripción funcional} \\
\midrule
1 & Autenticación y Perfil & Inicio de sesión, registro de conductores, recuperación de credenciales por correo electrónico, autenticación federada con Google OAuth2 y administración de perfiles y vehículos frecuentes. \\
2 & Portal del Conductor & Visualización gráfica de disponibilidad en mapa 2D, reserva de espacios, extensión de tiempo, historial de uso, calificaciones y pagos en línea mediante pasarela Stripe. \\
3 & Módulo de Garita & Escaneo y validación de código de reserva/QR, registro rápido de vehículos sin reserva (\textit{walk-in}), cálculo automatizado de tarifas al minuto, liquidación de salidas y solicitud digital de pago. \\
4 & Dashboard Administrativo & Indicadores clave de rendimiento (KPIs), gráficas interactivas de ingresos y ocupación en tiempo real, registro de actividad reciente y alertas operativas del sistema. \\
5 & Configuración y Mapa 2D & Herramienta visual e interactiva Drag \& Drop para diseñar la distribución física de cajones, definir tarifas dinámicas por tipo de vehículo, administrar al personal y consultar bitácoras de auditoría. \\
\bottomrule
\end{tabular}
\\ \vspace{0.2cm}
\textit{Nota.} Describe la función de cada uno de los cinco módulos que componen la plataforma ParkSys.
\end{table}

\subsubsection{Usuarios involucrados en el sistema}

\begin{table}[hbt!]
\caption{Usuarios involucrados en el sistema}
\label{tab:usuarios}
\begin{tabular}{@{}p{3cm} p{12cm}@{}}
\toprule
\textbf{Rol} & \textbf{Funciones principales} \\
\midrule
Conductor & Consulta disponibilidad de espacios en tiempo real, realiza reservas anticipadas, extiende su tiempo de estadía, efectúa pagos digitales en línea con tarjeta y califica la atención recibida. \\
Trabajador (Garita) & Controla físicamente las barreras de entrada y salida del estacionamiento, valida reservas existentes, registra ingresos no programados (\textit{walk-in}), cobra en efectivo o solicita pago digital y emite comprobantes. \\
Administrador & Supervisa el rendimiento financiero del negocio, configura tarifas por tipo de vehículo y turnos, modela la distribución geométrica del mapa 2D, gestiona las cuentas del personal y audita las operaciones del sistema. \\
\bottomrule
\end{tabular}
\\ \vspace{0.2cm}
\textit{Nota.} Resume las funciones y alcances operativos del conductor, el trabajador y el administrador.
\end{table}

\subsubsection{Problema que resuelve dentro de la organización}

\begin{itemize}
    \item \textbf{Elimina la ineficiencia operativa:} reemplaza los boletos físicos de papel y las libretas de control por registros electrónicos instantáneos en base de datos.
    \item \textbf{Erradica fugas de capital:} automatiza la liquidación económica mediante un algoritmo de cobro por minuto efectivo, impidiendo la manipulación deliberada o negligente de importes.
    \item \textbf{Mejora el control de personal:} cada ingreso, salida o cobro queda asociado de forma indeleble al identificador único del trabajador en turno.
    \item \textbf{Experiencia de usuario moderna (Contactless):} brinda al conductor la facultad de asegurar un espacio desde su dispositivo móvil, reduciendo tiempos de espera y emisión de gases contaminantes en la vía pública.
    \item \textbf{Visibilidad gerencial basada en datos:} facilita la toma de decisiones estratégicas mediante paneles analíticos que reflejan horas de mayor demanda, ingresos proyectados e índice de cumplimiento de reservas.
\end{itemize}

\begin{figure}[hbt!]
\caption{Diagrama general de la arquitectura del sistema ParkSys}
\label{fig:arquitectura_general}
% \includegraphics[width=\textwidth]{IMAGENES/figura_1.png}
\vspace{2cm} % Espacio reservado para la imagen
\begin{center}
\textit{[Espacio reservado para la Figura 1: Diagrama general de la arquitectura del sistema ParkSys]}
\end{center}
\vspace{2cm}
\textit{Nota.} Muestra la relación e interoperabilidad entre el Portal del Conductor, el Módulo de Garita, el Dashboard Administrativo, la API RESTful de FastAPI y la Base de Datos central PostgreSQL.
\end{figure}

\subsection{1.2. CONFIGURACIÓN DEL ENTORNO DE DESARROLLO}

Para el desarrollo de ParkSys se implementó un entorno de trabajo full-stack moderno, modular y completamente contenerizado, garantizando una rigurosa separación de entornos y dependencias.

\begin{table}[hbt!]
\caption{Lenguajes de programación empleados}
\label{tab:lenguajes}
\begin{tabular}{@{}p{2.5cm} p{2cm} p{2cm} p{8cm}@{}}
\toprule
\textbf{Componente} & \textbf{Lenguaje} & \textbf{Versión} & \textbf{Características y uso} \\
\midrule
Backend & Python & 3.11+ & Tipado estático con typing, concurrencia asíncrona nativa con asyncio y alto rendimiento de procesamiento. \\
Frontend & TypeScript & 5.5 & Superset de JavaScript con tipado estricto, interfaces DTO compiladas y soporte avanzado para decoradores de Angular. \\
\bottomrule
\end{tabular}
\\ \vspace{0.2cm}
\textit{Nota.} Detalla el lenguaje, la versión y la justificación técnica empleada en el backend y en el frontend.
\end{table}

\begin{table}[hbt!]
\caption{Frameworks principales}
\label{tab:frameworks}
\begin{tabular}{@{}p{2cm} p{3cm} p{10cm}@{}}
\toprule
\textbf{Capa} & \textbf{Framework} & \textbf{Justificación de uso} \\
\midrule
Backend & FastAPI 0.115 & Rendimiento a la par de NodeJS y Go gracias a Starlette y Pydantic v2; soporte nativo de asincronía (async/await), generación automática de documentación interactiva OpenAPI (Swagger UI) y validación declarativa de datos. \\
Frontend & Angular 18.2 & Arquitectura moderna basada en Standalone Components (sin sobrecarga de NgModules), reactividad granular mediante Angular Signals, excelente soporte empresarial para pruebas y navegación optimizada con Lazy Loading. \\
\bottomrule
\end{tabular}
\\ \vspace{0.2cm}
\textit{Nota.} Justifica la elección técnica de FastAPI en el backend y Angular en el frontend.
\end{table}

\begin{table}[hbt!]
\caption{Gestor de base de datos y acceso a datos}
\label{tab:bd}
\begin{tabular}{@{}p{3cm} p{3cm} p{9cm}@{}}
\toprule
\textbf{Herramienta} & \textbf{Versión} & \textbf{Función dentro de ParkSys} \\
\midrule
PostgreSQL & 16 (Alpine Linux) & Motor de base de datos relacional principal, seleccionado por su alta concurrencia, soporte nativo de transacciones ACID y tipos de datos avanzados (UUID, Timestamp with Timezone, Enums). \\
SQLAlchemy & 2.0 (modo asyncpg) & ORM en modo asíncrono no bloqueante que gestiona las transacciones y conexiones a la base de datos sin congelar el hilo principal de FastAPI. \\
Alembic & 1.13 & Sistema de control de versiones y migraciones evolutivas del esquema DDL de la base de datos. \\
\bottomrule
\end{tabular}
\\ \vspace{0.2cm}
\textit{Nota.} Detalla la función de PostgreSQL, SQLAlchemy y Alembic dentro del entorno de ParkSys.
\end{table}

\subsection{1.3. DISEÑO DE LA BASE DE DATOS}

La arquitectura de datos de ParkSys fue diseñada bajo un modelo relacional estrictamente normalizado en la Tercera Forma Normal (3FN), utilizando PostgreSQL 16. El diseño prioriza la integridad referencial, el aislamiento de transacciones para concurrencia masiva y la trazabilidad inmutable de eventos.

\subsubsection{Modelo Entidad-Relación (MER)}

\begin{table}[hbt!]
\caption{Descripción general de entidades y su propósito en el sistema}
\label{tab:entidades_general}
\begin{tabular}{@{}p{2.5cm} p{3.5cm} p{8.5cm}@{}}
\toprule
\textbf{Categoría} & \textbf{Entidad} & \textbf{Propósito principal} \\
\midrule
Núcleo Central & \texttt{users}, \texttt{spaces} & Almacena los perfiles con RBAC y el catálogo de cajones del estacionamiento con sus coordenadas 2D. \\
Transaccional & \texttt{reservations}, \texttt{vehicle\_logs}, \texttt{payments} & Maneja la vida útil de un servicio: reserva anticipada, ingreso físico, cálculo de duración y liquidación. \\
Configuración & \texttt{rates}, \texttt{parking\_info} & Parametriza las tarifas dinámicas y las políticas del establecimiento (tolerancia, horarios). \\
Legal y Soporte & \texttt{terms}, \texttt{reviews}, \texttt{audit\_log} & Controla versiones legales, calificaciones del servicio y la bitácora inmutable de acciones. \\
\bottomrule
\end{tabular}
\\ \vspace{0.2cm}
\textit{Nota.} Describe las 13 tablas principales que componen el ecosistema de datos del proyecto agrupadas por categoría funcional.
\end{table}

\begin{figure}[hbt!]
\caption{Modelo de Entidad-Relación (MER) lógico del sistema ParkSys}
\label{fig:mer}
\vspace{2cm} % Espacio reservado para la imagen
\begin{center}
\textit{[Espacio reservado para la Figura 2: Diagrama MER de ParkSys]}
\end{center}
\vspace{2cm}
\textit{Nota.} Representa la cardinalidad y relaciones entre las tablas transaccionales, el catálogo de usuarios y las métricas de auditoría.
\end{figure}

\subsection{1.4. DICCIONARIO DE DATOS}

El diccionario documenta la estructura de persistencia en PostgreSQL, generada a partir de los modelos asíncronos declarativos de SQLAlchemy 2.0.

\subsubsection{BLOQUE A: NÚCLEO DE USUARIOS Y CONTROL DE ESPACIOS}

\textbf{a) Tabla users — Archivo fuente: models/user.py.} Entidad central del sistema que administra el control de acceso (RBAC) y la identidad criptográfica de conductores, trabajadores y administradores.

\begin{table}[hbt!]
\caption{Diccionario de la tabla \texttt{users}}
\label{tab:dict_users}
\begin{tabular}{@{}p{2cm} p{2cm} p{1.5cm} p{1cm} p{1.5cm} p{6cm}@{}}
\toprule
\textbf{Campo} & \textbf{Tipo SQL} & \textbf{Longitud} & \textbf{Nulo} & \textbf{Clave} & \textbf{Descripción funcional} \\
\midrule
id & INTEGER & 4 bytes & No & PK & Identificador único del usuario. \\
nombre & VARCHAR & 100 & No & - & Nombre completo. \\
correo & VARCHAR & 100 & No & UK, Index & Medio principal de login y notificaciones. \\
password\_hash & VARCHAR & 255 & No & - & Contraseña cifrada en Argon2id. \\
rol & ENUM & - & No & - & Nivel de acceso: conductor, trabajador, admin. \\
created\_at & TIMESTAMPTZ & 8 bytes & No & - & Sello temporal de creación (UTC). \\
activo & BOOLEAN & 1 byte & No & - & Control de suspensión de cuenta (Soft-delete). \\
\bottomrule
\end{tabular}
\\ \vspace{0.2cm}
\textit{Nota.} Administra el control de acceso y la identidad criptográfica de los usuarios del sistema.
\end{table}

\textbf{b) Tabla spaces — Archivo fuente: models/space.py.} Representación digital y lógica de los cajones de estacionamiento.

\begin{table}[hbt!]
\caption{Diccionario de la tabla \texttt{spaces}}
\label{tab:dict_spaces}
\begin{tabular}{@{}p{2cm} p{2cm} p{1.5cm} p{1cm} p{1.5cm} p{6cm}@{}}
\toprule
\textbf{Campo} & \textbf{Tipo SQL} & \textbf{Longitud} & \textbf{Nulo} & \textbf{Clave} & \textbf{Descripción funcional} \\
\midrule
id & INTEGER & 4 bytes & No & PK & Identificador único del espacio. \\
codigo & VARCHAR & 10 & No & UK, Index & Nombre comercial (ej. A-01). \\
tipo\_vehiculo & ENUM & - & No & - & Restricción física (auto, moto). \\
estado & ENUM & - & No & - & Máquina de estados: disponible, ocupado, reservado. \\
reservable & BOOLEAN & 1 byte & No & - & Bandera para excluir espacios. \\
pos\_x & FLOAT & 8 bytes & No & - & Coordenada horizontal en plano 2D. \\
pos\_y & FLOAT & 8 bytes & No & - & Coordenada vertical en plano 2D. \\
\bottomrule
\end{tabular}
\\ \vspace{0.2cm}
\textit{Nota.} Define la geometría 2D, el estado operativo y las dimensiones de los espacios.
\end{table}

\subsubsection{BLOQUE B: TRANSACCIONALIDAD Y COBROS}

\textbf{c) Tabla reservations — Archivo fuente: models/reservation.py.} Registra la intención de ocupación a futuro mediante autoservicio.

\begin{table}[hbt!]
\caption{Diccionario de la tabla \texttt{reservations}}
\label{tab:dict_reservations}
\begin{tabular}{@{}p{2cm} p{2cm} p{1.5cm} p{1cm} p{1.5cm} p{6cm}@{}}
\toprule
\textbf{Campo} & \textbf{Tipo SQL} & \textbf{Longitud} & \textbf{Nulo} & \textbf{Clave} & \textbf{Descripción funcional} \\
\midrule
id & INTEGER & 4 bytes & No & PK & Identificador único. \\
usuario\_id & INTEGER & 4 bytes & No & FK, Index & Llave foránea hacia \texttt{users.id}. \\
espacio\_id & INTEGER & 4 bytes & No & FK, Index & Llave foránea hacia \texttt{spaces.id}. \\
placa & VARCHAR & 20 & No & - & Matrícula del vehículo esperado. \\
codigo\_reserva & VARCHAR & 12 & No & UK, Index & Token aleatorio de 8 dígitos alfanuméricos. \\
fecha\_reserva & TIMESTAMPTZ & 8 bytes & No & - & Creación de la reserva. \\
hora\_estimada & TIMESTAMPTZ & 8 bytes & No & - & Hora límite de tolerancia. \\
estado & ENUM & - & No & - & Estados: pendiente, confirmada, cancelada, vencida. \\
\bottomrule
\end{tabular}
\\ \vspace{0.2cm}
\textit{Nota.} Registra la intención de ocupación a futuro y la política de tolerancia.
\end{table}

\textbf{d) Tabla vehicle\_logs — Archivo fuente: models/vehicle\_log.py.} Tabla central de control de estadías efectivas (Walk-in y confirmación de reservas).

\begin{table}[hbt!]
\caption{Diccionario de la tabla \texttt{vehicle\_logs}}
\label{tab:dict_vehicle_logs}
\begin{tabular}{@{}p{2cm} p{2cm} p{1.5cm} p{1cm} p{1.5cm} p{6cm}@{}}
\toprule
\textbf{Campo} & \textbf{Tipo SQL} & \textbf{Longitud} & \textbf{Nulo} & \textbf{Clave} & \textbf{Descripción funcional} \\
\midrule
id & INTEGER & 4 bytes & No & PK & Identificador del registro físico. \\
placa & VARCHAR & 20 & No & Index & Matrícula registrada en garita. \\
espacio\_id & INTEGER & 4 bytes & No & FK, Index & Cajón asignado. \\
reserva\_id & INTEGER & 4 bytes & Sí & FK, Index & Nulo si es cliente walk-in. \\
fecha\_ingreso & TIMESTAMPTZ & 8 bytes & No & - & Apertura de barrera de ingreso. \\
fecha\_salida & TIMESTAMPTZ & 8 bytes & Sí & - & Apertura de barrera de salida. \\
estado & ENUM & - & No & - & Estados: en\_curso, finalizado. \\
\bottomrule
\end{tabular}
\\ \vspace{0.2cm}
\textit{Nota.} Tabla central de control de estadías efectivas que mide la ocupación física real.
\end{table}

\textbf{e) Tabla rates — Archivo fuente: models/rate.py.} Tarifario dinámico configurable por el administrador.

\begin{table}[hbt!]
\caption{Diccionario de la tabla \texttt{rates}}
\label{tab:dict_rates}
\begin{tabular}{@{}p{2cm} p{2cm} p{1.5cm} p{1cm} p{1.5cm} p{6cm}@{}}
\toprule
\textbf{Campo} & \textbf{Tipo SQL} & \textbf{Longitud} & \textbf{Nulo} & \textbf{Clave} & \textbf{Descripción funcional} \\
\midrule
id & INTEGER & 4 bytes & No & PK & Identificador de tarifa. \\
tipo\_vehiculo & ENUM & - & No & UK & auto, moto, camioneta. \\
precio\_hora & NUMERIC & 10,2 & No & - & Tarifa para la primera hora. \\
precio\_fraccion & NUMERIC & 10,2 & No & - & Costo por minuto adicional calculado matemáticamente. \\
vigente & BOOLEAN & 1 byte & No & - & Indica si esta tarifa es aplicable actualmente. \\
\bottomrule
\end{tabular}
\\ \vspace{0.2cm}
\textit{Nota.} Estructura el tarifario dinámico para el cálculo de cobros.
\end{table}

\textbf{f) Tabla payments — Archivo fuente: models/payment.py.} Registro inmutable de ingresos monetarios.

\begin{table}[hbt!]
\caption{Diccionario de la tabla \texttt{payments}}
\label{tab:dict_payments}
\begin{tabular}{@{}p{2cm} p{2cm} p{1.5cm} p{1cm} p{1.5cm} p{6cm}@{}}
\toprule
\textbf{Campo} & \textbf{Tipo SQL} & \textbf{Longitud} & \textbf{Nulo} & \textbf{Clave} & \textbf{Descripción funcional} \\
\midrule
id & INTEGER & 4 bytes & No & PK & Identificador del comprobante. \\
vehicle\_log\_id & INTEGER & 4 bytes & Sí & FK, Index & Estadía asociada. \\
reserva\_id & INTEGER & 4 bytes & Sí & FK, Index & Reserva asociada (pagos por penalidad). \\
monto & NUMERIC & 10,2 & No & - & Liquidación monetaria (soles). \\
metodo\_pago & ENUM & - & No & - & efectivo, tarjeta\_online, yape. \\
estado & ENUM & - & No & - & pendiente, completado, reembolsado. \\
fecha\_pago & TIMESTAMPTZ & 8 bytes & No & - & Momento de la transacción. \\
\bottomrule
\end{tabular}
\\ \vspace{0.2cm}
\textit{Nota.} Representa el registro inmutable de los ingresos monetarios.
\end{table}

\subsubsection{BLOQUE C: MARCO LEGAL Y AUDITORÍA INMUTABLE}

Debido a su naturaleza complementaria, las tablas restantes se resumen como entidades de soporte institucional y cumplimiento normativo.

\begin{table}[hbt!]
\caption{Diccionario resumen de las entidades de soporte y auditoría}
\label{tab:dict_soporte}
\begin{tabular}{@{}p{3cm} p{11.5cm}@{}}
\toprule
\textbf{Entidad} & \textbf{Propósito y descripción funcional principal} \\
\midrule
\texttt{parking\_info} & Registra metadatos institucionales: nombre, dirección, coordenadas GPS, horarios (apertura y cierre L-V, Sab, Dom) y minutos de tolerancia por No-Show. \\
\texttt{reviews} & Almacena comentarios y valoraciones (1 a 5 estrellas) emitidas por los conductores tras la culminación del servicio. Permite respuesta del administrador y moderación (visibilidad). \\
\texttt{notifications} & Almacena mensajes directos y alertas emitidas a la bandeja personal del usuario (cobros, liberaciones, advertencias). \\
\texttt{terms} & Administra las versiones publicadas (ej. "v1.0") de los términos de servicio, políticas de privacidad y condiciones de reserva. \\
\texttt{terms\_acceptance} & Registra de forma inmutable el consentimiento legal explícito otorgado por cada usuario (\texttt{usuario\_id}) sobre una versión particular de los términos (\texttt{terms\_id}), con fecha exacta. \\
\texttt{audit\_log} & Pista de auditoría inalterable (Audit Trail) para fiscalizar acciones administrativas sensibles (modificación de tarifas, anulación de reservas, eliminación de usuarios). Carece de operación DELETE. \\
\bottomrule
\end{tabular}
\\ \vspace{0.2cm}
\textit{Nota.} Sintetiza el propósito de las tablas de configuración, mensajería, legales y de trazabilidad (Tablas 20 a 26 del informe original).
\end{table}

\subsection{1.5. DESARROLLO DEL FRONTEND}

La interfaz de usuario de ParkSys fue construida priorizando la fluidez, la velocidad de respuesta y una estética visual corporativa moderna, incorporando principios de Glassmorphism y microinteracciones fluidas.

\subsubsection{Arquitectura de navegación y rendimiento}

\begin{itemize}
    \item \textbf{Lazy Loading:} los módulos de admin, worker, conductor y auth son empaquetados en chunks independientes de JavaScript, descargándose bajo demanda únicamente cuando el usuario navega a sus rutas específicas.
    \item \textbf{Componentes Standalone:} se eliminó la sobrecarga estructural de NgModule, reduciendo el tamaño del bundle final y facilitando el aislamiento de dependencias.
\end{itemize}

\subsubsection{Diseño de experiencia de usuario (UX/UI)}

\begin{table}[hbt!]
\caption{Diseño de experiencia de usuario (UX/UI)}
\label{tab:ux_ui}
\begin{tabular}{@{}p{4cm} p{10cm}@{}}
\toprule
\textbf{Elemento de diseño} & \textbf{Especificación técnica implementada} \\
\midrule
Paleta de colores corporativa & Azul medianoche profundo (\#021024), azul marino oscuro (\#052659), azul acero (\#5483B3) y azul hielo (\#C1E8FF), combinados con fondos oscuros. \\
Estilo visual & Glassmorphism elegante: bordes delgados semitransparentes (border-white/10), desenfoque de fondo (backdrop-blur-md) y sombras sutiles. \\
Animaciones & Transiciones fluidas en entrada de listas (fade-in-up), elevación táctil (hover:translate-y-1), ondas de clic (ripples) y pulsos cromáticos. \\
Semántica cromática & Verde esmeralda (libres), rojo carmesí (ocupados), ámbar cálido (reservados) y gris grafito (mantenimiento). \\
\bottomrule
\end{tabular}
\\ \vspace{0.2cm}
\textit{Nota.} Resume la identidad visual, la paleta cromática y los estándares estéticos aplicados en la interfaz.
\end{table}

\subsubsection{Gestión del estado reactivo}

A diferencia de arquitecturas Angular tradicionales basadas exclusivamente en RxJS, ParkSys integra activamente \textbf{Angular Signals} para la reactividad de interfaz. RxJS se mantiene exclusivamente en la capa de infraestructura para manejar flujos asíncronos desacoplados (peticiones HTTP y WebSockets).

\begin{table}[hbt!]
\caption{Cuadro comparativo del enfoque reactivo}
\label{tab:enfoque_reactivo}
\begin{tabular}{@{}p{3cm} p{5.5cm} p{6cm}@{}}
\toprule
\textbf{Enfoque reactivo} & \textbf{Caso de uso principal} & \textbf{Ventaja técnica observada} \\
\midrule
RxJS (Observables) & Peticiones HTTP al backend, reconexión de sockets y manejo de streams temporizados. & Manejo robusto de flujos asíncronos y cancelación de solicitudes con \texttt{switchMap}. \\
Angular Signals & Estado visual síncrono: cajón seleccionado, modales activos, contadores. & Cero sobrecarga de detección de cambios (Zoneless-ready) y reactividad granular. \\
\bottomrule
\end{tabular}
\\ \vspace{0.2cm}
\textit{Nota.} Compara la complementariedad técnica entre RxJS y Angular Signals dentro del frontend.
\end{table}

\begin{figure}[hbt!]
\caption{Arquitectura de Lazy Loading del frontend}
\label{fig:lazy_loading}
\vspace{2cm} % Espacio reservado para la imagen
\begin{center}
\textit{[Espacio reservado para la Figura 11: Arquitectura Lazy Loading]}
\end{center}
\vspace{2cm}
\textit{Nota.} Representa la división y carga diferida de los módulos de Angular.
\end{figure}

\subsection{1.6. INTEGRACIÓN DE MÓDULOS}

La comunicación integral del ecosistema se logra a través del protocolo HTTPS REST (para operaciones transaccionales) y WSS (WebSockets sobre TLS para sincronización push instantánea).

\begin{table}[hbt!]
\caption{Flujo integrado del proceso de reserva}
\label{tab:flujo_reserva}
\begin{tabular}{@{}p{0.5cm} p{3.5cm} p{2cm} p{8cm}@{}}
\toprule
\textbf{Paso} & \textbf{Emisor $\rightarrow$ Receptor} & \textbf{Protocolo} & \textbf{Acción técnica ejecutada} \\
\midrule
1 & Conductor $\rightarrow$ Frontend & Local UI & Selecciona espacio en mapa 2D; se computa estado reactivo con Signals. \\
2 & Frontend $\rightarrow$ Backend & POST REST & Envía payload a \texttt{/reservations} con \texttt{espacio\_id}, placa y hora. \\
3 & Backend $\rightarrow$ PostgreSQL & SQLAlchemy & Valida concurrencia y ejecuta el INSERT de la reserva. \\
4 & Backend $\rightarrow$ PostgreSQL & SQLAlchemy & Actualiza el estado de la tabla \texttt{spaces} a 'reservado'. \\
5 & Backend $\rightarrow$ WebSockets & Push Frame & Emite broadcast \texttt{\{"event": "refresh"\}} a todos los clientes. \\
6 & WebSockets $\rightarrow$ Garita & Local UI & Panel del operador conmuta el cajón a color ámbar inmediatamente. \\
\bottomrule
\end{tabular}
\\ \vspace{0.2cm}
\textit{Nota.} Detalla paso a paso la coreografía entre el cliente web, la API RESTful y el motor de base de datos.
\end{table}

\begin{figure}[hbt!]
\caption{Secuencia completa del flujo de reserva}
\label{fig:flujo_reserva_secuencia}
\vspace{2cm} % Espacio reservado para la imagen
\begin{center}
\textit{[Espacio reservado para la Figura 12: Diagrama de secuencia]}
\end{center}
\vspace{2cm}
\textit{Nota.} Ilustra el ciclo completo desde la selección del cajón hasta la actualización en caseta.
\end{figure}

\subsection{1.7. VALIDACIÓN DE DATOS Y MANEJO DE ERRORES}

La robustez y resiliencia de ParkSys se garantiza mediante un modelo defensivo de doble anillo: validación síncrona en cliente y validación semántica estricta en servidor.

\subsubsection{Validación en el Frontend y Backend}

En el frontend se emplea \texttt{ReactiveFormsModule} con validadores síncronos (obligatoriedad, correo bien formado, expresiones regulares para placas). En el backend, cada payload JSON es interceptado por modelos Pydantic v2. Si una solicitud contiene tipos incorrectos, FastAPI responde automáticamente con \texttt{HTTP 422 Unprocessable Entity}.

\begin{table}[hbt!]
\caption{Cuadro de códigos de error gestionados}
\label{tab:codigos_error}
\begin{tabular}{@{}p{1.5cm} p{5.5cm} p{3.5cm} p{4cm}@{}}
\toprule
\textbf{Código HTTP} & \textbf{Escenario de activación} & \textbf{Componente de origen} & \textbf{Respuesta emitida al cliente} \\
\midrule
400 & Solicitud incorrecta (monto inferior a tarifa mínima). & Controladores & JSON con mensaje explicativo. \\
401 & Token JWT ausente, alterado o expirado. & Dependencies & Redirección forzada hacia \texttt{/login}. \\
403 & Rol del usuario insuficiente (RBAC). & \texttt{require\_role()} & Notificación de acceso denegado. \\
404 & Identificador solicitado no existe. & SQLAlchemy query & Notificación en pantalla (no hallado). \\
409 & Conflicto de estado (doble reserva concurrente). & Lógica de reservas & Mensaje de colisión de concurrencia. \\
422 & Estructura JSON malformada o discordante. & Pydantic Validator & Desglose de restricciones infringidas. \\
500 & Falla no controlada del servidor. & Middleware Global & Traza en logs y error neutralizado. \\
\bottomrule
\end{tabular}
\\ \vspace{0.2cm}
\textit{Nota.} Enumera los códigos de estado HTTP manejados formalmente por la plataforma y su tratamiento.
\end{table}

\subsection{1.8. IMPLEMENTACIÓN DE LA SEGURIDAD Y AUTENTICACIÓN}

ParkSys implementa un esquema de seguridad sin estado (Stateless), escalable y alineado con los lineamientos de OWASP Top 10. Las contraseñas se cifran mediante el algoritmo de derivación de claves Argon2id, y la autenticación se gestiona mediante tokens JWT simétricos (HS256) con ciclo de vida segmentado (Access Token de 15 minutos y Refresh Token de 7 días).

\begin{table}[hbt!]
\caption{Autorización basada en roles (RBAC)}
\label{tab:rbac}
\begin{tabular}{@{}p{2.5cm} p{12cm}@{}}
\toprule
\textbf{Rol} & \textbf{Nivel de acceso otorgado en el sistema} \\
\midrule
admin & Acceso total: lectura y escritura en tarifas, reportes financieros, editor 2D, administración de personal y consulta de auditoría inmutable. \\
trabajador & Acceso operativo a garita: verificación de reservas, registro de clientes \textit{walk-in}, visualización de mapa y terminal de cobros. \\
conductor & Acceso de autoservicio: exploración de mapa, apartado de cajones, pagos online con tarjeta, historial personal y gestión de matrículas. \\
\bottomrule
\end{tabular}
\\ \vspace{0.2cm}
\textit{Nota.} Describe formalmente la matriz de privilegios y niveles de acceso autorizados.
\end{table}

\begin{table}[hbt!]
\caption{Cuadro resumen de las capas de seguridad implementadas}
\label{tab:seguridad_capas}
\begin{tabular}{@{}p{3.5cm} p{7cm} p{4cm}@{}}
\toprule
\textbf{Capa de seguridad} & \textbf{Mecanismo técnico} & \textbf{Herramienta / Estándar} \\
\midrule
Persistencia de Claves & Hashing criptográfico unidireccional con sal & Argon2id (passlib) \\
Autenticación & Tokens criptográficos sin estado con expiración & JSON Web Tokens (HS256) \\
Identidad Federada & Verificación asimétrica de tokens de terceros & Google OAuth2 \\
Autorización (RBAC) & Control de acceso por roles inyectable & \texttt{require\_role()} (FastAPI) \\
Rutas en Frontend & Guards reactivos de navegación cliente & authGuard y roleGuard \\
Protección de Dominio & Control de orígenes y cabeceras anti-XSS & CORS Whitelist, nosniff \\
Trazabilidad Legal & Registro inmutable de eventos y consentimientos & \texttt{audit\_log}, \texttt{terms\_acceptance} \\
\bottomrule
\end{tabular}
\\ \vspace{0.2cm}
\textit{Nota.} Sintetiza los mecanismos defensivos en las distintas capas del sistema.
\end{table}

\begin{figure}[hbt!]
\caption{Flujo completo de autenticación y autorización}
\label{fig:flujo_auth}
\vspace{2cm} % Espacio reservado para la imagen
\begin{center}
\textit{[Espacio reservado para la Figura 14 y 15: Seguridad y Autenticación]}
\end{center}
\vspace{2cm}
\textit{Nota.} Muestra la secuencia desde el inicio de sesión hasta la verificación de permisos en endpoints protegidos.
\end{figure}


\clearpage
\section{CAPÍTULO 2: PRUEBAS, OPTIMIZACIÓN Y DESPLIEGUE}

\subsection{2.1. PRUEBAS UNITARIAS, DE INTEGRACIÓN Y DE ACEPTACIÓN}




\subsubsection{Estrategia general de pruebas}




Para dar estricto cumplimiento a las directrices del Módulo Formativo IV (Implementación de Pruebas de Software), el aseguramiento de la calidad de ParkSys se ejecutó bajo la metodología de la Pirámide de Pruebas de Software, segmentando la verificación en tres niveles formales. La estrategia adoptada consistió en realizar una auditoría endpoint por endpoint tanto en el Backend (API REST con FastAPI) como en el Frontend (Angular 18), registrando para cada funcionalidad la tecnología utilizada, los errores detectados en la primera versión, las correcciones aplicadas en la segunda versión y las recomendaciones pendientes para futuras iteraciones.

\begin{figure}[H]
\caption{Pirámide de pruebas de software aplicada a ParkSys}
\centering
\includegraphics[width=\textwidth,height=9cm,keepaspectratio]{IMAGENES/figuras/figura_16.png}
\vspace{0.2cm}
{\raggedright \textit{Nota.} Representa los tres niveles de la pirámide de pruebas —pruebas unitarias en la base, pruebas de integración en el nivel intermedio y pruebas de aceptación (UAT) en la cúspide— aplicados sobre los módulos billing.py, security.py, schemas.py, la API REST, SQLAlchemy, PostgreSQL y WebSockets. \par}
\end{figure}

\subsubsection{Herramientas y tecnologías utilizadas}




\begin{table}[hbt!]
\caption{Herramientas y tecnologías empleadas en las pruebas de ParkSys}
\label{tab:herramientas_pruebas}
\begin{tabular}{@{}p{3cm} p{1.5cm} p{8cm} p{2.5cm}@{}}
\toprule
\textbf{Herramienta} & \textbf{Versión} & \textbf{Propósito} & \textbf{Capa} \\
\midrule
pytest & 8.x & Framework principal de pruebas en el backend Python & Backend \\
pytest-asyncio & 0.23+ & Soporte para pruebas asíncronas con async/await & Backend \\
httpx & 0.27 & Cliente HTTP asíncrono para pruebas de endpoints FastAPI & Backend \\
aiosqlite & 0.20+ & Driver SQLite asíncrono para base de datos en memoria & Backend \\
SQLAlchemy AsyncSession & 2.0.35 & ORM asíncrono para crear fixtures y verificar persistencia & Backend \\
Jasmine & 5.x & Framework de pruebas para componentes Angular 18 & Frontend \\
Karma & 6.x & Test runner para ejecutar pruebas Jasmine en navegador & Frontend \\
Angular Signals & 18.x & Detección de cambios reactiva para validar re-renders & Frontend \\
RxJS & 7.x & Pruebas de flujos reactivos (WebSocket, Observables) & Frontend \\
\bottomrule
\end{tabular}
\\ \vspace{0.2cm}
\textit{Nota.} Enumera las herramientas de prueba utilizadas en cada capa del sistema, con su versión y propósito técnico.
\end{table}

\subsubsection{Cuadros consolidados de evaluación endpoint por endpoint}




\begin{table}[hbt!]
\caption{Pruebas unitarias y de integración en endpoints — Backend (FastAPI / Python)}
\label{tab:pruebas_backend}
\begin{tabular}{@{}p{2.5cm} p{1.8cm} p{2.2cm} p{3.0cm} p{3.5cm} p{3cm}@{}}
\toprule
\textbf{Endpoint} & \textbf{Tipo prueba} & \textbf{Tech} & \textbf{V1: tiempo/errores} & \textbf{V2: correcciones} & \textbf{Recomendaciones} \\
\midrule
POST /auth/login & Unit / Integr. & pytest, httpx & 120 ms. Error HTTP 500 si usuario no existía. & Manejo HTTP 401. Tiempo: 85 ms. & Rate limiting (5/min). \\
POST /auth/register & Unitaria & pytest, Pydantic & 95 ms. Aceptaba contraseñas cortas y correos inválidos. & Validación estricta con Pydantic. Tiempo: 80 ms. & Verificación correo por token. \\
POST /reservations & Integración & pytest, SQLAlch. & 200 ms. Condición de carrera en reservas concurrentes. & Validación dentro de transacción atómica. Tiempo: 110 ms. & SELECT FOR UPDATE. \\
GET /spaces & Unitaria & pytest, httpx & 150 ms. Consultas N+1 por reserva. & joinedload() e índice B-Tree. Tiempo: 42 ms. & Caché Redis (TTL 5s). \\
POST /payments & Integración & pytest, Stripe & 300 ms. Timeout y errores en montos menores a S/2.00. & Reintentos y validación de monto mínimo. Tiempo: 250 ms. & Webhooks Stripe en 2do plano. \\
POST /vehicles & Unitaria & pytest, asyncio & 140 ms. Aceptaba placas nulas/vacías. & Validación regex estricta (Pydantic). Tiempo: 90 ms. & OCR en garita. \\
PATCH /vehicles/../checkout & Integración & pytest, SQLAlch. & 190 ms. Cálculo incorrecto de duración a medianoche. & math.ceil y timezone.utc consistentes. Tiempo: 18 ms. & Alerta por >24 horas de estadía. \\
GET /reports/pdf/id & Integración & pytest, ReportLab & 350 ms. Desfase UTC de 5 horas. & zoneinfo.ZoneInfo("America/Lima"). Tiempo: 280 ms. & Worker Celery+RabbitMQ para PDF. \\
\bottomrule
\end{tabular}
\\ \vspace{0.2cm}
\textit{Nota.} Consolida la evolución de las pruebas backend en el formato Versión 1 $\rightarrow$ Versión 2 $\rightarrow$ Recomendaciones.
\end{table}

\begin{table}[hbt!]
\caption{Pruebas de consumo de endpoints — Frontend (Angular 18 / TypeScript)}
\label{tab:pruebas_frontend}
\begin{tabular}{@{}p{3.5cm} p{2.0cm} p{2cm} p{2.5cm} p{3cm} p{2.5cm}@{}}
\toprule
\textbf{Componente} & \textbf{Tipo prueba} & \textbf{Tech} & \textbf{V1: tiempo/errores} & \textbf{V2: correcciones} & \textbf{Recomendaciones} \\
\midrule
AuthService $\rightarrow$ login & Unitaria & Jasmine, Karma & 180 ms. UI bloqueada sin mensaje ante credencial errónea. & JwtInterceptor y MatSnackBar. Redirección rol. & Pruebas E2E (Cypress). \\
RegisterComponent & Unitaria & Jasmine, Karma & 200 ms. Sin error en tiempo real. & Validaciones FormGroup reactivas. & Debounce validación correo. \\
NewReservationComponent & Unitaria & Jasmine, Karma & 250 ms. Mapa no refrescaba tras reserva. & Angular Signals y evento WebSocket. & Animaciones (fade-in ámbar). \\
MapComponent & Integración & Jasmine, RxJS & 300 ms. Polling HTTP de alto consumo. & WebSocketService; fin de polling. & Canvas/WebGL para $>1000$ espacios. \\
GaritaPanelComponent & Integración & Jasmine, Karma & 200 ms. Desconexión silenciosa WS. & Reconexión automática (3s) y alerta. & Modo offline (PWA). \\
PaymentComponent & Integración & Jasmine, HttpClient & 280 ms. Doble clic sin spinner. & Deshabilitación y spinner; validación monto. & Integrar Yape/Plin. \\
CheckoutComponent & Integración & Jasmine, Karma & 220 ms. PDF no descargaba auto. & blob y window.open(). & Impresora POS térmica. \\
\bottomrule
\end{tabular}
\\ \vspace{0.2cm}
\textit{Nota.} Consolida la evolución de las pruebas frontend en el formato Versión 1 $\rightarrow$ Versión 2 $\rightarrow$ Recomendaciones.
\end{table}

\begin{table}[hbt!]
\caption{Reporte de errores generales y optimización de rendimiento por módulo}
\label{tab:errores_generales}
\begin{tabular}{@{}p{2.5cm} p{2cm} p{2cm} p{3.5cm} p{3cm} p{3cm}@{}}
\toprule
\textbf{Módulo/Capa} & \textbf{Tipo análisis} & \textbf{Tech} & \textbf{V1: problemas} & \textbf{V2: mejoras} & \textbf{Recomendaciones} \\
\midrule
Base de datos & Rendimiento & PostgreSQL & Consultas lentas por placa, queries O(N). & Índices B-Tree en 10 campos. O(log N). & Particionar vehicle\_logs. \\
Facturación & Error & ReportLab & Desfase UTC en PDF. & zoneinfo America/Lima en backend/front. & PDFs asíncronos. \\
Frontend SPA & Rendimiento & Angular 18 & FCP de 3.8s en móviles. & Standalone Components y Lazy Loading. FCP 1.1s. & Habilitar SSR (Angular Univ.). \\
API REST & Error & FastAPI & CORS bloqueado en Railway. & Regex allow\_origin\_regex. & CSP y restricción CORS. \\
WebSockets & Error & FastAPI+Ang. & Desconexión silenciosa. & Reconexión automática. & Ping/pong heartbeat. \\
Pasarela & Error & Stripe API & HTTP 400 por $<$ S/ 2.00. & Derivación a efectivo. & MercadoPago alternativo. \\
\bottomrule
\end{tabular}
\\ \vspace{0.2cm}
\textit{Nota.} Sintetiza los hallazgos y mejoras de rendimiento y estabilidad a nivel de sistema completo.
\end{table}

\subsubsection{2.1.1. Pruebas unitarias (Unit Testing)}




Las pruebas unitarias verificaron de forma aislada las funciones puras y algoritmos críticos del sistema, sin dependencia de bases de datos externas ni servidores HTTP. Se ejecutaron 30 pruebas unitarias distribuidas en tres componentes fundamentales del backend.

\textbf{Archivo fuente: backend/tests/test\_unit\_billing.py} \\
\textbf{a) Motor de facturación (billing.py) — 14 pruebas}\\
Se verificó que el algoritmo calculase exactamente los minutos efectivos a partir de la diferencia entre fecha\_ingreso y fecha\_salida, redondeando hacia arriba (math.ceil).

\begin{verbatim}
# UT-04 — Redondeo hacia arriba (ceil)
def test_ut04_redondeo_hacia_arriba(self):
    ingreso = datetime(2026, 8, 15, 10, 0, 0, tzinfo=timezone.utc)
    salida = ingreso + timedelta(seconds=90)
    result = calculate_duration_minutes(ingreso, salida)
    assert result == 2
\end{verbatim}

\textbf{b) Seguridad criptográfica (security.py) — 10 pruebas}\\
Se verificó la integridad del hash Argon2id y validación JWT.

\begin{verbatim}
# UT-15 — Hash Argon2 nunca es texto plano
def test_ut15_hash_no_es_texto_plano(self):
    plain = "admin123"
    hashed = hash_password(plain)
    assert hashed != plain
    assert "$argon2" in hashed
\end{verbatim}

\textbf{c) Validación de esquemas Pydantic (schemas.py) — 6 pruebas}
\begin{verbatim}
# UT-26 — Password menor a 8 caracteres
def test_ut26_registro_password_corto(self):
    with pytest.raises(ValidationError):
        RegisterRequest(nombre="Test", correo="t@t.com", password="123")
\end{verbatim}

\begin{table}[hbt!]
\caption{Matriz completa de pruebas unitarias (30 pruebas)}
\label{tab:matriz_ut}
\begin{tabular}{@{}p{2.5cm} p{3.5cm} p{4cm} p{3.5cm} p{2cm}@{}}
\toprule
\textbf{ID} & \textbf{Componente} & \textbf{Entrada} & \textbf{Salida esperada} & \textbf{Resultado} \\
\midrule
UT-01 a UT-08 & billing.py / calculate\_dur & Casos de duración exactos/fracc. & Minutos con redondeo ceil & Conforme \\
UT-09 a UT-14 & billing.py / calculate\_chr & Tipo veh., minutos, tarifa & Monto total calculado & Conforme \\
UT-15 a UT-19 & security.py / hash, verify & Contraseñas prueba & Hash Argon2, validación ok & Conforme \\
UT-20 a UT-24 & security.py / JWT & Payloads acceso/refresco & Tokens gen/decod/rech & Conforme \\
UT-25 a UT-30 & schemas.py & Datos válidos/inválidos & ValidationError o ok & Conforme \\
\bottomrule
\end{tabular}
\\ \vspace{0.2cm}
\textit{Nota.} Resume la matriz completa de 30 pruebas unitarias.
\end{table}

\subsubsection{2.1.2. Pruebas de integración (Integration Testing)}




Se ejecutaron 12 pruebas de integración.

\begin{verbatim}
# IT-02 — Crear reserva cambia espacio a 'reservado'
async def test_it02_crear_reserva_cambia_espacio_a_reservado(...):
    assert sample_space.estado == SpaceStatus.disponible
    # Insert logic...
    assert space_db.estado == SpaceStatus.reservado
\end{verbatim}

\begin{table}[hbt!]
\caption{Matriz completa de pruebas de integración (12 pruebas)}
\label{tab:matriz_it}
\begin{tabular}{@{}p{2cm} p{4cm} p{4cm} p{3cm} p{2cm}@{}}
\toprule
\textbf{ID} & \textbf{Escenario} & \textbf{Procedimiento} & \textbf{Criterio de éxito} & \textbf{Resultado} \\
\midrule
IT-01a/b/c & Autenticación endpoints & Petición sin token/inválido & HTTP 401/403/200 & Conforme \\
IT-02 & Reserva $\rightarrow$ reservado & Crear reserva cajón disp. & Espacio actualizado & Conforme \\
IT-03 & Doble reserva placa & Segunda reserva placa activa & Validación impide dup & Conforme \\
IT-04, 11-12 & Búsquedas indexadas & Consulta código, placa & Búsqueda eficiente & Conforme \\
IT-05 & Cancelación libera espacio & Cancelar reserva pendiente & Espacio disponible & Conforme \\
IT-06, 07 & Transición garita & Walk-in y checkout & Estado cambia ocup/disp & Conforme \\
IT-08 & Cajón bloqueado & reservable=False & HTTP 409 Rechazo & Conforme \\
IT-09, 10 & Bitácora vehicle\_logs & Ingreso/salida walk-in & Log persistido con mins & Conforme \\
\bottomrule
\end{tabular}
\\ \vspace{0.2cm}
\textit{Nota.} Resume la matriz completa de 12 pruebas de integración.
\end{table}

\subsubsection{2.1.3. Pruebas de aceptación (UAT)}

Diseñadas bajo el formato formal Gherkin (Dado / Cuando / Entonces), simulando los flujos completos que los usuarios ejecutarían en producción en el local ParkSys Central Ayacucho. Se ejecutaron 6 pruebas de aceptación.

\textbf{Archivo fuente:} \texttt{backend/tests/test\_acceptance\_flows.py}

\textbf{Ejemplo completo — UAT-03: Cobro exacto al minuto}
\begin{verbatim}
class TestUAT03CobroExacto:
 """
 DADO: Un vehículo permaneció exactamente 1h 15min en el estacionamiento.
 CUANDO: El trabajador presiona "Dar Salida" en el panel.
 ENTONCES: Se liquida hora base + fracción proporcional exacta = S/ 6.25.
 """
 async def test_uat03_cobro_exacto_con_fraccion(self, db_session):
     fecha_ingreso = datetime(2026, 8, 20, 10, 0, 0, tzinfo=timezone.utc)
     fecha_salida = datetime(2026, 8, 20, 11, 15, 0, tzinfo=timezone.utc)
     charge = await calculate_charge(db_session, VehicleType.auto,
                                     fecha_ingreso, fecha_salida)
     assert charge["total"] == 6.25
     assert charge["full_hours"] == 1
     assert charge["remaining_minutes"] == 15
\end{verbatim}





Se ejecutaron 6 pruebas de aceptación (Gherkin).

\begin{table}[hbt!]
\caption{Matriz de pruebas de aceptación (6 pruebas)}
\label{tab:matriz_uat}
\begin{tabular}{@{}p{4cm} p{3.5cm} p{3.5cm} p{4.5cm}@{}}
\toprule
\textbf{Escenario} & \textbf{Condición previa (Dado)} & \textbf{Acción (Cuando)} & \textbf{Resultado (Entonces)} \\
\midrule
UAT-01: Reserva exitosa & Conductor autenticado & Fija fecha/hora y reserva & Código 8 dig, notificación \\
UAT-02: Admisión walk-in & Operador en garita & Digita placa, confirma & Autoasignación, espacio ocupado \\
UAT-03: Cobro exacto al min & Vehículo permaneció 75min & Pulsa "Dar Salida" & 75 min * tarifa = S/6.25, recibo \\
UAT-04: WebSocket push & Operador con WS activo & Conductor crea reserva web & Evento refresco en $<$ 200ms \\
UAT-05: Cancelación reserva & Conductor reserva B-03 & Pulsa "Cancelar" & Cajón disp, reserva cancelada \\
UAT-06: Expiración auto & Reserva vencida & APScheduler ejecuta & Estado vencida, cajón liberado \\
\bottomrule
\end{tabular}
\\ \vspace{0.2cm}
\textit{Nota.} Sintetiza los seis flujos de aceptación validados en formato Gherkin.
\end{table}

\begin{figure}[H]
\caption{Resumen del total de pruebas ejecutadas por nivel}
\centering
\includegraphics[width=\textwidth,height=9cm,keepaspectratio]{IMAGENES/figuras/figura_17.png}
\vspace{0.2cm}
{\raggedright \textit{Nota.} Gráfico de barras que muestra la distribución de las 48 pruebas automatizadas ejecutadas (30 unitarias, 12 de integración y 6 de aceptación), con una tasa de éxito del 100 por ciento en todas ellas. \par}
\end{figure}

\subsection{2.2. REPORTE DE ERRORES}

Durante las etapas de desarrollo, integración y despliegue del sistema se registraron incidencias técnicas que fueron debidamente categorizadas, diagnosticadas y subsanadas. Cada error fue documentado con su identificador, componente afectado, descripción detallada de la falla, nivel de impacto, causa raíz técnica y la solución aplicada.




\begin{table}[hbt!]
\caption{Registro formal de errores identificados y resueltos}
\label{tab:registro_errores_v2}
\begin{tabular}{@{}p{1.2cm} p{2cm} p{4cm} p{1.8cm} p{2.5cm} p{3.5cm}@{}}
\toprule
\textbf{ID} & \textbf{Capa} & \textbf{Falla} & \textbf{Impacto} & \textbf{Causa} & \textbf{Solución} \\
\midrule
ERR-01 & Base de Datos & Condición de carrera: reservas duplicadas simultáneas. & Alto & Sin bloqueo transaccional concurrente. & SELECT ... WHERE estado='disponible' (expire\_on\_commit=False). \\
ERR-02 & WebSockets & Desconexión silenciosa en inactividad. & Medio & Sin reconexión, micro-cortes WiFi. & Reconexión (3s) en Angular, limpieza de sockets backend. \\
ERR-03 & Zonas Horarias & Desfase de 5h UTC en PDF. & Medio & Servidor UTC, sede UTC-5. & zoneinfo.ZoneInfo("America/Lima"). \\
ERR-04 & Stripe API & HTTP 400 por cobro $<$ S/ 2.00. & Medio & Mínimo de Stripe 0.50 USD. & Validación backend: derivar a efectivo. \\
ERR-05 & CORS Deploy & CORS bloqueado en subdominios Railway. & Alto & ALLOWED\_ORIGINS estático. & Regex allow\_origin\_regex en middleware CORS. \\
\bottomrule
\end{tabular}
\\ \vspace{0.2cm}
\textit{Nota.} Registra los cinco errores críticos y de impacto medio.
\end{table}

\begin{figure}[H]
\caption{Flujo de diagnóstico y resolución de errores (ERR-01 a ERR-05)}
\centering
\includegraphics[width=\textwidth,height=9cm,keepaspectratio]{IMAGENES/figuras/figura_18.png}
\vspace{0.2cm}
{\raggedright \textit{Nota.} Diagrama de flujo que ilustra el proceso seguido para cada incidencia. \par}
\end{figure}

\begin{table}[hbt!]
\caption{Distribución de errores por nivel de impacto}
\label{tab:distribucion_errores}
\begin{tabular}{@{}p{6cm} p{3cm} p{6cm}@{}}
\toprule
\textbf{Nivel de impacto} & \textbf{Cantidad} & \textbf{Identificadores} \\
\midrule
Alto (bloqueante para producción) & 2 & ERR-01, ERR-05 \\
Medio (afecta experiencia/precisión) & 3 & ERR-02, ERR-03, ERR-04 \\
Bajo (mejoras de UX sin bloqueo) & 2 & Hallazgos UAT en sección 2.4 \\
\bottomrule
\end{tabular}
\\ \vspace{0.2cm}
\textit{Nota.} Clasifica los errores documentados según su severidad.
\end{table}

\subsection{2.3. OPTIMIZACIÓN DEL RENDIMIENTO}

Para asegurar que ParkSys mantenga tiempos de respuesta mínimos bajo condiciones de alta concurrencia en horas punta, se ejecutaron optimizaciones de ingeniería en dos capas: Frontend y Backend.




\subsubsection{2.3.1. Frontend: Lazy Loading y Standalone Components}

Angular 18 permite definir cada componente como Standalone (standalone: true), eliminando la necesidad de NgModules tradicionales. ParkSys aprovecha esta capacidad junto con Lazy Loading para reducir drásticamente el peso inicial de la aplicación. Los módulos de administrador, trabajador y conductor se cargan solo cuando el usuario navega a su sección correspondiente mediante \texttt{loadComponent()} en el archivo \texttt{app.routes.ts}.




\begin{table}[hbt!]
\caption{Componentes con Lazy Loading implementado}
\label{tab:lazy_loading_v2}
\begin{tabular}{@{}p{3cm} p{9cm} p{3cm}@{}}
\toprule
\textbf{Módulo} & \textbf{Componentes Lazy-Loaded} & \textbf{Ahorro est.} \\
\midrule
Administrador & Financials, Personnel, Settings, Support, Audit, Clients, AdminReviews, TransactionDetail & $\sim$ 180 KB \\
Conductor & SecuritySettings, NotificationSettings, EditProfile, Receipt, PendingPayments, PaymentSuccess & $\sim$ 120 KB \\
Público & ResetPassword, Terms, PublicLegal, PublicSupport, PublicPaymentSuccess & $\sim$ 60 KB \\
\bottomrule
\end{tabular}
\\ \vspace{0.2cm}
\textit{Nota.} Detalla la distribución de componentes cargados de forma diferida.
\end{table}

\begin{figure}[H]
\caption{Comparativa de First Contentful Paint antes y después de la optimización}
\centering
\includegraphics[width=\textwidth,height=9cm,keepaspectratio]{IMAGENES/figuras/figura_19.png}
\vspace{0.2cm}
{\raggedright \textit{Nota.} Gráfico comparativo que muestra la reducción del FCP de 3.8 segundos a 1.1 segundos. \par}
\end{figure}

\subsubsection{2.3.2. Backend y base de datos: asyncpg e indexación}

\textbf{Driver asíncrono asyncpg}

El uso del driver nativo asyncpg con SQLAlchemy 2.0 Async permite que FastAPI atienda múltiples peticiones concurrentes sin bloquear el hilo principal esperando respuestas de disco:

\begin{verbatim}
# backend/app/db/session.py
engine = create_async_engine(
 settings.database_url, # postgresql+asyncpg://...
 echo=settings.environment == "development",
 pool_pre_ping=True,
)

async_session_factory = async_sessionmaker(
 engine,
 class_=AsyncSession,
 expire_on_commit=False,
)
\end{verbatim}

\textbf{Resultado:} el servidor soporta más de 800 peticiones por segundo en benchmarks locales con uvicorn y el pool de conexiones de asyncpg.

\textbf{Índices B-Tree en campos de búsqueda frecuente}

Se crearon 10 índices explícitos en las columnas más consultadas, reduciendo la complejidad de búsqueda de O(N) a O(log N).




\begin{table}[hbt!]
\caption{Campos indexados en el sistema}
\label{tab:indices_bd_v2}
\begin{tabular}{@{}p{3cm} p{3.5cm} p{3.5cm} p{5cm}@{}}
\toprule
\textbf{Tabla} & \textbf{Campo} & \textbf{Tipo de índice} & \textbf{Uso principal} \\
\midrule
users & correo & UNIQUE & Login y registro \\
spaces & codigo & UNIQUE & Búsqueda de cajón \\
reservations & codigo\_reserva & UNIQUE & Búsqueda en garita \\
reservations & usuario\_id & B-Tree FK & Listado por conductor \\
reservations & espacio\_id & B-Tree FK & Verificación de disponibilidad \\
vehicle\_logs & placa & B-Tree & Búsqueda por placa \\
vehicle\_logs & reserva\_id & B-Tree FK & Join reservations \\
vehicle\_logs & espacio\_id & B-Tree FK & Historial por espacio \\
payments & vehicle\_log\_id & B-Tree FK & Asociación pago-registro \\
payments & reserva\_id & B-Tree FK & Asociación pago-reserva \\
\bottomrule
\end{tabular}
\\ \vspace{0.2cm}
\textit{Nota.} Enumera los diez índices creados sobre columnas de mayor frecuencia.
\end{table}

\begin{table}[hbt!]
\caption{Métricas de optimización y benchmarks}
\label{tab:metricas_optimizacion}
\begin{tabular}{@{}p{5cm} p{3cm} p{3cm} p{4cm}@{}}
\toprule
\textbf{Indicador} & \textbf{Valor base} & \textbf{Valor opt.} & \textbf{Mejora} \\
\midrule
First Contentful Paint (FCP) & 3.8 s & 1.1 s & 71\% más rápido \\
Latencia /spaces/availability & 280 ms & 42 ms & 85\% de reducción \\
Tiempo de cálculo de checkout & 190 ms & 18 ms & 90\% de reducción \\
Consumo de RAM en reposo & 180 MB & 78 MB & 56\% más liviano \\
Throughput del servidor & 320 req/s & 800+ req/s & 150\% de incremento \\
\bottomrule
\end{tabular}
\\ \vspace{0.2cm}
\textit{Nota.} Sintetiza las métricas de rendimiento antes y después.
\end{table}

\begin{figure}[H]
\caption{Arquitectura de indexación y acceso asíncrono a la base de datos}
\centering
\includegraphics[width=\textwidth,height=9cm,keepaspectratio]{IMAGENES/figuras/figura_20.png}
\vspace{0.2cm}
{\raggedright \textit{Nota.} Diagrama que ilustra el flujo con asyncpg e índices B-Tree. \par}
\end{figure}

\begin{figure}[H]
\caption{Comparativa de throughput del servidor antes y después de la optimización}
\centering
\includegraphics[width=\textwidth,height=9cm,keepaspectratio]{IMAGENES/figuras/figura_21.png}
\vspace{0.2cm}
{\raggedright \textit{Nota.} Gráfico de barras que compara el throughput del servidor. \par}
\end{figure}

\subsection{2.4. PRUEBAS DE USUARIO (UAT)}

Las pruebas de usuario se ejecutaron en un entorno controlado con usuarios reales representando los tres roles del sistema: Administrador, Trabajador de garita y Conductor. El objetivo fue validar que los flujos de trabajo del sistema se ejecutaran de manera intuitiva, eficiente y libre de errores desde la perspectiva del usuario final, y no exclusivamente desde la del desarrollador.




\begin{table}[hbt!]
\caption{Parámetros de ejecución de las pruebas de usuario}
\label{tab:parametros_uat}
\begin{tabular}{@{}p{4cm} p{11cm}@{}}
\toprule
\textbf{Parámetro} & \textbf{Detalle} \\
\midrule
Fecha de ejecución & Agosto del 2026 \\
Lugar & Sede ParkSys Central — Ayacucho \\
Participantes & Tres usuarios (un administrador, un trabajador y un conductor) \\
Dispositivos empleados & PC de escritorio (garita), laptop (admin), smartphone Android \\
Entorno de prueba & Producción (Railway), conectado a PostgreSQL remoto \\
Metodología aplicada & Observación directa, cuestionario post-uso y registro de incidencias \\
\bottomrule
\end{tabular}
\\ \vspace{0.2cm}
\textit{Nota.} Describe las condiciones de ejecución de las pruebas.
\end{table}

\begin{figure}[H]
\caption{Evidencia fotográfica de la sesión de pruebas de usuario en la sede ParkSys Central}
\centering
\includegraphics[width=\textwidth,height=9cm,keepaspectratio]{IMAGENES/figuras/figura_22.png}
\vspace{0.2cm}
{\raggedright \textit{Nota.} Registra el momento de la ejecución de las pruebas con usuarios operando el sistema. \par}
\end{figure}

\begin{table}[hbt!]
\caption{Escenarios de prueba — Rol Conductor}
\label{tab:uat_conductor}
\begin{tabular}{@{}p{0.5cm} p{4.5cm} p{4.5cm} p{1.8cm} p{3.7cm}@{}}
\toprule
\textbf{N.°} & \textbf{Tarea} & \textbf{Criterio de éxito} & \textbf{Result.} & \textbf{Observaciones} \\
\midrule
1 & Registro correo/pass & Cuenta creada e inicio & Cumplido & 45 segundos \\
2 & Agregar vehículo & Listado en "Mis Vehículos" & Cumplido & Sin dificultad \\
3 & Crear reserva mapa 2D & Código y cajón en ámbar & Cumplido & Muy intuitivo \\
4 & Consultar estado & Estado "pendiente" & Cumplido & - \\
5 & Cancelar reserva & Estado "cancelada" & Cumplido & Sugirió confirmación UI \\
6 & Pagar Stripe online & Checkout exitoso/recibo & Cumplido & Redirect de 3 seg \\
\bottomrule
\end{tabular}
\\ \vspace{0.2cm}
\textit{Nota.} Detalla las tareas ejecutadas por el conductor.
\end{table}

\begin{table}[hbt!]
\caption{Escenarios de prueba — Rol Trabajador de Garita}
\label{tab:uat_trabajador}
\begin{tabular}{@{}p{0.5cm} p{3.7cm} p{4cm} p{1.8cm} p{4.5cm}@{}}
\toprule
\textbf{N.°} & \textbf{Tarea} & \textbf{Criterio de éxito} & \textbf{Result.} & \textbf{Observaciones} \\
\midrule
1 & Iniciar sesión panel & Mapa en tiempo real & Cumplido & Carga en 1.2s \\
2 & Ingreso walk-in & Espacio/hora registrada & Cumplido & $<$ 10 seg \\
3 & Dar salida (checkout) & Monto mostrado, PDF & Cumplido & PDF automático \\
4 & Verificación mapa WS & Espacio cambia color auto & Cumplido & $<$ 200 ms via WebSocket \\
5 & Buscar por placa & Resultado rápido & Cumplido & 40 ms indexado \\
\bottomrule
\end{tabular}
\\ \vspace{0.2cm}
\textit{Nota.} Describe las tareas del trabajador de garita.
\end{table}

\begin{table}[hbt!]
\caption{Escenarios de prueba — Rol Administrador}
\label{tab:uat_admin}
\begin{tabular}{@{}p{0.5cm} p{3.7cm} p{4cm} p{1.8cm} p{4.5cm}@{}}
\toprule
\textbf{N.°} & \textbf{Tarea} & \textbf{Criterio de éxito} & \textbf{Result.} & \textbf{Observaciones} \\
\midrule
1 & Acceder dashboard KPIs & Indicadores visibles & Cumplido & Vista consolidada OK \\
2 & Gestionar espacios & CRUD cajones & Cumplido & - \\
3 & Gestionar personal & CRUD trabajadores & Cumplido & - \\
4 & Configurar tarifas & Tarifas reflejadas & Cumplido & Verificó cobro moto \\
5 & Reportes financieros & Filtrados por fechas & Cumplido & - \\
6 & Revisar auditoría & Registro temporal OK & Cumplido & - \\
\bottomrule
\end{tabular}
\\ \vspace{0.2cm}
\textit{Nota.} Detalla las tareas administrativas del administrador.
\end{table}

\begin{table}[hbt!]
\caption{Resultados consolidados de las pruebas de usuario}
\label{tab:resultados_uat_v2}
\begin{tabular}{@{}p{8cm} p{7cm}@{}}
\toprule
\textbf{Métrica} & \textbf{Resultado} \\
\midrule
Tasa de completitud de tareas & 100\% (17 de 17 tareas completadas) \\
Errores críticos durante pruebas & 0 \\
Errores menores (sugerencias UX) & 2 (Confirmación cancelación, auto-mayúsculas) \\
Satisfacción general (escala 1-5) & 4.7 \\
Tiempo promedio aprendizaje & Menos de 3 minutos por rol \\
\bottomrule
\end{tabular}
\\ \vspace{0.2cm}
\textit{Nota.} Sintetiza los indicadores cuantitativos de UAT.
\end{table}

\subsection{2.5. DOCUMENTACIÓN TÉCNICA Y MANUAL DE USUARIO}

\subsubsection{2.5.1. Documentación Técnica (Dirigida al Desarrollador)}

La documentación técnica de ParkSys se encuentra distribuida en los siguientes artefactos del repositorio del proyecto:




\begin{table}[hbt!]
\caption{Documentos técnicos del repositorio}
\label{tab:documentos_tecnicos}
\begin{tabular}{@{}p{4cm} p{4cm} p{7.5cm}@{}}
\toprule
\textbf{Documento} & \textbf{Ubicación} & \textbf{Contenido} \\
\midrule
README.md & ./README.md & Descripción, stack, instalación \\
CAPITULO\_1.md & ./CAPITULO\_1.md & Análisis, modelos, arquitectura API \\
CAPITULO\_2\_PRUEBAS & ./CAP...\_PRUEBAS.md & Unit/Integration/UAT, errores \\
DESCRIPCION.md & ./DESCRIPCION.md & Requisitos funcionales \\
.env.example & ./backend/.env... & Plantilla de entorno \\
pyproject.toml & ./backend/py... & Config pytest asíncrono \\
docker-compose.yml & ./docker-compose... & Orquestación local (PG + FastAPI) \\
Dockerfile & ./backend/Docker... & Imagen Python Uvicorn \\
Swagger OpenAPI & /docs (FastAPI) & Documentación interactiva de API \\
\bottomrule
\end{tabular}
\\ \vspace{0.2cm}
\textit{Nota.} Enumera artefactos documentales.
\end{table}

\begin{figure}[H]
\caption{Arquitectura del sistema (Resumen técnico de capas)}
\centering
\includegraphics[width=\textwidth,height=9cm,keepaspectratio]{IMAGENES/figuras/figura_23.png}
\vspace{0.2cm}
{\raggedright \textit{Nota.} Ilustra la disposición en capas del sistema: el frontend en Angular 18 con sus tres módulos por rol, el backend en FastAPI con los routers, la capa de servicios y el ORM asíncrono, y la base de datos PostgreSQL 16 con sus índices B-Tree. \par}
\end{figure}

\begin{table}[hbt!]
\caption{Endpoints de la API REST (15 routers registrados)}
\label{tab:endpoints_api}
\begin{tabular}{@{}p{3.5cm} p{3.5cm} p{8cm}@{}}
\toprule
\textbf{Prefijo} & \textbf{Router} & \textbf{Descripción} \\
\midrule
/api/v1/auth & auth.py & Login, registro, token refresh, OAuth \\
/api/v1/spaces & spaces.py & CRUD cajones, 2D mapa \\
/api/v1/reservations & reservations.py & Reservas, cancelación \\
/api/v1/vehicles & vehicles.py & Walk-in, checkout \\
/api/v1/payments & payments.py & Stripe checkout, transacciones \\
/api/v1/reviews & reviews.py & Reseñas conductores \\
/api/v1/notifications & notifications.py & Alertas, mensajes \\
/api/v1/reports & reports.py & Reportes PDF (ReportLab) \\
/api/v1/settings & settings.py & Tarifas, horarios, tolerancias \\
/api/v1/workers & workers.py & Personal garita \\
/api/v1/history & history.py & Historial log vehicular \\
/api/v1/terms & terms.py & Términos legales \\
/api/v1/audit & audit.py & Bitácora inmutable \\
/api/v1/users & users.py & CRUD perfiles \\
/api/v1/ws & websockets.py & Real-time WebSockets \\
\bottomrule
\end{tabular}
\\ \vspace{0.2cm}
\textit{Nota.} Conjunto de controladores REST de la API.
\end{table}

\begin{table}[hbt!]
\caption{Modelos de base de datos referenciados en la documentación}
\label{tab:modelos_bd}
\begin{tabular}{@{}p{3cm} p{3cm} p{9cm}@{}}
\toprule
\textbf{Modelo} & \textbf{Tabla} & \textbf{Campos principales} \\
\midrule
User & users & id, nombre, correo, password\_hash, rol \\
Space & spaces & id, codigo, tipo\_vehiculo, estado, reservable... \\
Reservation & reservations & id, usuario\_id, espacio\_id, placa, estado \\
VehicleLog & vehicle\_logs & id, placa, ingreso, salida \\
Payment & payments & id, vehicle\_log\_id, monto, estado \\
Rate & rates & id, tipo\_vehiculo, precio\_hora \\
ParkingInfo & parking\_info & id, horarios, nombre, direccion \\
\bottomrule
\end{tabular}
\\ \vspace{0.2cm}
\textit{Nota.} Referencia rápida de ORMs.
\end{table}

\subsubsection{2.5.2. Manual de Usuario (Dirigido al Usuario Final)}

El manual de usuario fue diseñado para los tres perfiles del sistema, con instrucciones secuenciales para cada flujo de trabajo principal.

\begin{table}[hbt!]
\caption{Manual del Conductor}
\label{tab:manual_conductor_v2}
\begin{tabular}{@{}p{0.5cm} p{10.5cm} p{4cm}@{}}
\toprule
\textbf{Paso} & \textbf{Acción} & \textbf{Pantalla} \\
\midrule
1 & Ingresar a la URL del sistema web & Landing Page \\
2 & Pulsar "Registrarse" y completar formulario & Registro \\
3 & Iniciar sesión con correo y contraseña & Login \\
4 & Ingresar a "Mis Vehículos" y agregar & Vehículos \\
5 & Ingresar a "Nueva Reserva" y seleccionar cajón 2D & Mapa 2D \\
6 & Fijar fecha/hora y reservar & Reserva \\
7 & Anotar código alfanumérico & Confirmación \\
8 & Consultar estado en "Mis Reservas" & Listado \\
9 & Pago Stripe online / Efectivo & Pago \\
10 & Descargar comprobante PDF & Historial \\
\bottomrule
\end{tabular}
\\ \vspace{0.2cm}
\textit{Nota.} Procedimiento paso a paso para el conductor.
\end{table}

\begin{table}[hbt!]
\caption{Manual del Trabajador de Garita}
\label{tab:manual_trabajador_v2}
\begin{tabular}{@{}p{0.5cm} p{10.5cm} p{4cm}@{}}
\toprule
\textbf{Paso} & \textbf{Acción} & \textbf{Pantalla} \\
\midrule
1 & Iniciar sesión & Login \\
2 & Visualizar panel en tiempo real & Panel Garita \\
3 & Ingreso con reserva (validar código 8 dígitos) & Búsqueda \\
4 & Ingreso walk-in (digitar placa, auto-asignar) & Form walk-in \\
5 & Dar salida por placa (Checkout autom.) & Checkout \\
6 & Registrar pago (efectivo/tarjeta) & Registro pago \\
7 & Comprobante PDF entregado & Comprobante \\
\bottomrule
\end{tabular}
\\ \vspace{0.2cm}
\textit{Nota.} Operativa en garita.
\end{table}

\begin{table}[hbt!]
\caption{Manual del Administrador}
\label{tab:manual_admin_v2}
\begin{tabular}{@{}p{0.5cm} p{10.5cm} p{4cm}@{}}
\toprule
\textbf{Paso} & \textbf{Acción} & \textbf{Pantalla} \\
\midrule
1 & Iniciar sesión admin & Login \\
2 & Visualizar KPIs en tiempo real & Dashboard \\
3 & CRUD de cajones & Gestor espacios \\
4 & CRUD personal trabajador & Personal \\
5 & Tarifas, horarios, tolerancias & Configuración \\
6 & Reportes PDF de caja & Finanzas \\
7 & Auditoría temporal de eventos & Auditoría \\
8 & Listado clientes & Clientes \\
\bottomrule
\end{tabular}
\\ \vspace{0.2cm}
\textit{Nota.} Funciones administrativas.
\end{table}

\subsection{2.6. DESPLIEGUE DEL SISTEMA}




\begin{figure}[H]
\caption{Arquitectura de despliegue de ParkSys (entorno de producción y entorno local)}
\centering
\includegraphics[width=\textwidth,height=9cm,keepaspectratio]{IMAGENES/figuras/figura_24.png}
\vspace{0.2cm}
{\raggedright \textit{Nota.} Ilustra la disposición de los servicios en producción y local. \par}
\end{figure}

\begin{table}[hbt!]
\caption{Configuración del servicio de base de datos (entorno local)}
\label{tab:config_bd_local}
\begin{tabular}{@{}p{4cm} p{11cm}@{}}
\toprule
\textbf{Parámetro} & \textbf{Valor} \\
\midrule
Imagen & postgres:16-alpine \\
Base de datos & parksys \\
Puerto expuesto & 5432 \\
Volumen persistente & postgres\_data \\
Healthcheck & pg\_isready, cada 5 segundos, 5 reintentos \\
\bottomrule
\end{tabular}
\\ \vspace{0.2cm}
\textit{Nota.} Contenedor PostgreSQL local.
\end{table}

\begin{table}[hbt!]
\caption{Configuración del servicio backend (entorno local)}
\label{tab:config_backend_local}
\begin{tabular}{@{}p{4cm} p{11cm}@{}}
\toprule
\textbf{Parámetro} & \textbf{Valor} \\
\midrule
Build & ./backend/Dockerfile \\
Imagen base & python:3.11-slim \\
Puerto & 8000 \\
Dependencia & Espera healthcheck PostgreSQL \\
Comando & alembic upgrade head \&\& uvicorn app.main:app \\
\bottomrule
\end{tabular}
\\ \vspace{0.2cm}
\textit{Nota.} Contenedor FastAPI local.
\end{table}

\begin{table}[hbt!]
\caption{Criterios de selección de la plataforma de despliegue (Railway)}
\label{tab:criterios_railway}
\begin{tabular}{@{}p{3.5cm} p{12cm}@{}}
\toprule
\textbf{Criterio} & \textbf{Justificación} \\
\midrule
Costo & Capa gratuita MVP \\
Facilidad & CI/CD automático con GitHub Webhook \\
PostgreSQL & Soporte nativo manejado con backups \\
Docker & Soporte nativo de Dockerfile \\
Certificados & HTTPS automático let's encrypt \\
Variables entorno & Panel unificado amigable \\
\bottomrule
\end{tabular}
\\ \vspace{0.2cm}
\textit{Nota.} Justifica Railway como PaaS productivo.
\end{table}

\begin{table}[hbt!]
\caption{Configuración del backend en Railway}
\label{tab:config_railway_backend}
\begin{tabular}{@{}p{4cm} p{11cm}@{}}
\toprule
\textbf{Parámetro} & \textbf{Valor} \\
\midrule
Directorio raíz & /backend \\
Comando Build & pip install -r requirements.txt \\
Comando Inicio & alembic upgrade head \&\& uvicorn ... --port \$PORT \\
Uso Dockerfile & Sí \\
\bottomrule
\end{tabular}
\\ \vspace{0.2cm}
\textit{Nota.} Backend productivo.
\end{table}

\begin{table}[hbt!]
\caption{Configuración del frontend en Railway}
\label{tab:config_railway_frontend}
\begin{tabular}{@{}p{4cm} p{11cm}@{}}
\toprule
\textbf{Parámetro} & \textbf{Valor} \\
\midrule
Directorio raíz & /frontend \\
Comando Build & npm install \&\& ng build --configuration production \\
Dir de salida & dist/frontend/browser \\
Tipo de servicio & Sitio estático SPA \\
\bottomrule
\end{tabular}
\\ \vspace{0.2cm}
\textit{Nota.} Frontend productivo SPA.
\end{table}

\begin{figure}[H]
\caption{Flujo del pipeline de despliegue continuo (CI/CD) de ParkSys}
\centering
\includegraphics[width=\textwidth,height=9cm,keepaspectratio]{IMAGENES/figuras/figura_25.png}
\vspace{0.2cm}
{\raggedright \textit{Nota.} Ilustra la secuencia desde GitHub hasta producción automática. \par}
\end{figure}

\begin{table}[hbt!]
\caption{Variables de entorno configuradas en el entorno de producción}
\label{tab:env_produccion}
\begin{tabular}{@{}p{6cm} p{9cm}@{}}
\toprule
\textbf{Variable} & \textbf{Descripción} \\
\midrule
DATABASE\_URL & Conexión PG (Railway inject) \\
JWT\_SECRET & Firma HS256 JWT \\
STRIPE\_SECRET\_KEY & Pasarela de pagos secreta \\
ALLOWED\_ORIGINS & Dominios CORS permitidos \\
SMTP\_USER & Credenciales correo alertas \\
\bottomrule
\end{tabular}
\\ \vspace{0.2cm}
\textit{Nota.} Sintetiza el .env productivo.
\end{table}

\subsection{2.7. ENTREGA Y PUESTA EN MARCHA}




\begin{table}[hbt!]
\caption{Checklist de entregables del proyecto ParkSys}
\label{tab:checklist_entregables}
\begin{tabular}{@{}p{1cm} p{4cm} p{7cm} p{2cm}@{}}
\toprule
\textbf{N.°} & \textbf{Entregable} & \textbf{Descripción} & \textbf{Estado} \\
\midrule
1 & Código fuente & Repositorio Git completo & Entregado \\
2 & Base de datos & PG 16 migrado + seed & Entregado \\
3 & Backend desplegado & FastAPI activo en Railway & Operativo \\
4 & Frontend desplegado & Angular 18 activo & Operativo \\
5 & Documentación & Manuales técnicos MD & Entregado \\
6 & Manual Usuario & Manuales por perfil & Entregado \\
7 & Seed Data & seed.py ejecutado & Ejecutado \\
8 & Credenciales & Cuentas admin/trab/cond & Entregado \\
9 & Suite de pruebas & 48 automatizadas (100\% ok) & Entregado \\
10 & Docker local & Config docker-compose & Entregado \\
\bottomrule
\end{tabular}
\\ \vspace{0.2cm}
\textit{Nota.} Consolida el estado de cumplimiento entregables.
\end{table}

\begin{table}[hbt!]
\caption{Fases del procedimiento de puesta en marcha}
\label{tab:fases_puesta_marcha}
\begin{tabular}{@{}p{0.5cm} p{9.5cm} p{2cm} p{2cm}@{}}
\toprule
\textbf{Fase} & \textbf{Actividad} & \textbf{Respon.} & \textbf{Estado} \\
\midrule
1 & Verificación infraestructura red garita & Técnico & Cumplido \\
2 & Despliegue producción Railway & Dev & Cumplido \\
3 & Migración DDL alembic & Dev & Cumplido \\
4 & Carga inicial seed.py & Dev & Cumplido \\
5 & Configuración tarifas/tolerancia admin & Admin & Cumplido \\
6 & Diseño mapa 2D interactivo & Admin & Cumplido \\
7 & Capacitación de garita (walk-in/checkout) & Dev/Trabaj & Cumplido \\
8 & Prueba completa real E2E & Todos & Cumplido \\
9 & Acta Aceptación & Admin & Cumplido \\
\bottomrule
\end{tabular}
\\ \vspace{0.2cm}
\textit{Nota.} Describe las 9 fases de puesta en marcha operativa.
\end{table}

\begin{figure}[H]
\caption{Diagrama de flujo del procedimiento de puesta en marcha del sistema ParkSys}
\centering
\includegraphics[width=\textwidth,height=9cm,keepaspectratio]{IMAGENES/figuras/figura_26.png}
\vspace{0.2cm}
{\raggedright \textit{Nota.} Ilustra la secuencia cronológica de puesta en marcha. \par}
\end{figure}

\begin{table}[hbt!]
\caption{Plan de contingencia post-entrega}
\label{tab:plan_contingencia}
\begin{tabular}{@{}p{4cm} p{8cm} p{3cm}@{}}
\toprule
\textbf{Riesgo} & \textbf{Acción} & \textbf{Responsable} \\
\midrule
Caída servidor Railway & Auto-restart; o levantar Docker local. & Dev \\
Pérdida internet garita & Log manual temporal (libreta), sync posterior. & Trabajador \\
Error Base de datos & Restaurar backup diario Railway automático. & Dev \\
Bug productivo & Hotfix git push $\rightarrow$ auto-deploy. & Dev \\
Aumento de tráfico & Scale up de memoria/CPU en Railway panel. & Dev \\
\bottomrule
\end{tabular}
\\ \vspace{0.2cm}
\textit{Nota.} Acciones previstas ante escenarios operativos.
\end{table}

\begin{table}[hbt!]
\caption{Acta de aceptación final del proyecto ParkSys}
\label{tab:acta_aceptacion}
\begin{tabular}{@{}p{4cm} p{11cm}@{}}
\toprule
\textbf{Campo} & \textbf{Detalle} \\
\midrule
Proyecto & ParkSys — Sistema de Gestión Estacionamiento \\
Cliente & Estacionamiento Central — Ayacucho \\
Fecha entrega & Agosto 2026 \\
Equipo & Aspur Yaranga, Condori Cotaquispe \\
Módulos entregados & Conductor, Garita, Admin, Stripe, WS, PDF, Aud \\
Pruebas ejecutadas & 48 automatizadas, 17 tareas reales \\
Tasa de éxito & 100\% ok \\
Conformidad & Aceptado sin observaciones \\
\bottomrule
\end{tabular}
\\ \vspace{0.2cm}
\textit{Nota.} Documenta formalmente la conformidad del cliente.
\end{table}


\clearpage
\section{CAPÍTULO 3: CONCLUSIONES Y RECOMENDACIONES}

\subsection{3.1. CONCLUSIONES}

\begin{enumerate}
    \item El desarrollo del proyecto ParkSys, durante las Experiencias Formativas en Situaciones Reales de Trabajo IV, permitió consolidar de manera práctica los conocimientos adquiridos en el módulo de Implementación de Pruebas de Software, evidenciando que la aplicación sistemática de pruebas unitarias, de integración y de aceptación (48 pruebas automatizadas con una tasa de éxito del 100 por ciento) constituye un factor determinante para garantizar la confiabilidad de un sistema antes de su despliegue en un entorno de producción real.
    \item La experiencia formativa permitió comprobar que el proceso de optimización del rendimiento —desde la identificación de cuellos de botella hasta la implementación de mejoras concretas, tales como la indexación de campos de alta consulta y la adopción de Lazy Loading en el frontend— exige no solo dominio técnico, sino también una metodología ordenada de diagnóstico, corrección y verificación, la cual fue aplicada de forma iterativa a lo largo del desarrollo del sistema.
    \item El trabajo colaborativo realizado durante el desarrollo, las pruebas y el despliegue de ParkSys permitió fortalecer competencias transversales relacionadas con la gestión de proyectos de software en equipo, la documentación técnica rigurosa y la comunicación efectiva con el usuario final, competencias que resultan indispensables para el desempeño profesional del futuro Ingeniero de Sistemas de Información.
\end{enumerate}

\subsection{3.2. RECOMENDACIONES}

\begin{enumerate}
    \item Se recomienda que la escuela profesional destine mayor tiempo curricular a la etapa de pruebas de software dentro del plan de Experiencias Formativas, dado que, durante el desarrollo del presente proyecto, se evidenció que la elaboración de una suite de pruebas completa (unitarias, de integración y de aceptación) demanda una carga de trabajo considerable que en ocasiones compite con los plazos establecidos para las demás etapas del proyecto.
    \item Se sugiere fortalecer, dentro de los cursos previos a las Experiencias Formativas, la enseñanza de herramientas de control de versiones y de despliegue continuo (CI/CD), ya que su dominio resultó fundamental para la entrega oportuna del sistema y podría reforzarse con mayor práctica guiada antes de enfrentar un proyecto real.
    \item Se recomienda a la coordinación académica establecer espacios de asesoría periódica más frecuentes durante la etapa de implementación de pruebas y despliegue, con la finalidad de absolver oportunamente las dudas técnicas que surgen durante estas fases y optimizar así los tiempos de ejecución del proyecto.
\end{enumerate}


\end{document}
